\subsection{AFFINE LINEAR VARIETIES}

DEFINITION 2. Given an affine space E, a subset V of E is called an affine linear variety (or simply a linear variety or an affine subset of E) if, for every family \((x_{t})_{t\in I}\) of points of V and every family \((\lambda_t)_{t\in I}\) of elements of K of finite support such that \(\sum_{t\in I}\lambda_t = 1\) , the barycentre \(\sum_{t\in I}\lambda_t x_t\) belongs to V.

It amounts to the same to say that the condition of Definition 2 holds for every finite family of points of V.

The empty set is a linear variety; every intersection of linear varieties is a linear variety.

Let V be a non-empty subset of E and \(a\) a point of V; the relation

\[
x - a = \sum_ {i = 1} ^ {n} \lambda_ {i} (x _ {i} - a)
\]

means that \(x\) is a barycentre \(\sum_{i=1}^{n} \lambda_i x_i + \left(1 - \sum_{i=1}^{n} \lambda_i\right)a\) of the family consisting of the \(x_i\) and \(a\) . Therefore:

PROPOSITION 2. For a non-empty subset V of an affine space E to be a linear variety, it is necessary and sufficient that V be a vector subspace for the vector space structure on E obtained by taking a point of V as origin.

In particular, the non-empty affine linear varieties of a vector space T (considered as an affine space) are just the translates of the vector subspaces of T; the vector subspaces of T are therefore the linear varieties containing 0.

Let V be a non-empty linear variety of the affine space E; the set of free vectors x - y, where x and y run through V, is a vector subspace D of the translation space T of E called the direction of V: for, if \(a \in V\) , then

\[
x - y = (x - a) - (y - a)
\]

and our assertion follows from Proposition 2. It is immediate that D operates faithfully and transitively on V, which therefore has canonically the structure of an affine space attached to D. By the dimension of the affine variety V, we mean the dimension of V with this affine space structure, that is the dimension of the vector space D. The linear varieties of dimension 0 are the points of E; those of dimension 1 (resp. 2) are called lines (resp. planes) of E.

Every vector \(\neq0\) belonging to the direction of a line is called a direction vector of this line; its components with respect to a basis of T form what is called a system of direction parameters of the line in question.

The codimension of a linear variety V in E is the codimension of its direction D in T; a linear variety of codimension 1 in E is called an (affine) hyperplane of E.

Two linear varieties with the same direction are called parallel; it amounts to the same to say that one is derived from the other by translation. If V is a linear variety in T (considered as an affine space), its direction is the linear variety parallel to V and containing 0.

PROPOSITION 3. Given a family \((a_{t})_{t\in I}\) of points of an affine space \(\mathbf{E}\) , the set \(\mathbf{V}\) of barycentres \(\sum_{i\in I}\lambda_i a_i\) (( \(\lambda_i\) ) of finite support, \(\sum_{i\in I}\lambda_i = 1\) ) is a linear variety of \(\mathbf{E}\) .

If the family \((a_{t})\) is empty, then \(\mathbf{V} = \varnothing\) because of the condition \(\sum_{i}\lambda_{i} = 1\) . It may therefore be assumed that the family \((a_{t})\) is non-empty and in this case the proposition is obvious, taking one of the \(a_{t}\) as origin in E.

The variety V is obviously the smallest linear variety containing the \(a_{t}\) ; it is said to be generated by the family \((a_{t})\) and this family is called a generating system of V.

In the notation of Proposition 3, assuming the family \((a_{\iota})\) is non-empty, for the expression for every point \(x \in V\) in the form \(x = \sum_{\iota} \lambda_{\iota} a_{\iota}\) to be unique, it is necessary and sufficient that, denoting an arbitrary index of I by \(\kappa\) , the family of vectors \(a_{\iota} - a_{\kappa}\) , where \(\iota\) runs through the set of indices \(\neq \kappa\) , be free in T. Then the family \((a_{\iota})_{\iota \in I}\) of points of E is said to be affinely free (or that its elements form an affinely free system, or are affinely independent) and that \(\lambda_{\iota}\) is the barycentric coordinate of \(x\) of index \(\iota\) with respect to the affinely free family \((a_{\iota})\) .

A family \((a_{i})_{i\in I}\) of points of E which is not affinely free is said to be affinely related.

PROPOSITION 4. For a non-empty family \((a_{t})_{t\in I}\) of points in an affine space \(E\) to be affinely related, it is necessary and sufficient that there exist a family \((\lambda_t)_{t\in I}\) of elements not all zero in \(K\) , of finite support, such that \(\sum_{i\in I}\lambda_i = 0\) and \(\sum_{i\in I}\lambda_i a_i = 0\) .

Given an index \(\kappa \in I\) , to say that the family of vectors \((a_{t} - a_{\kappa})\) , where \(t\) runs through the set of indices \(\neq \kappa\) , is related in \(T\) , means that there exists a family of scalars \((\lambda_{t})_{t \neq \kappa}\) not all zero such that \(\sum_{t \neq \kappa} \lambda_{t}(a_{t} - a_{\kappa}) = 0\) , which may also be written \(\sum_{t \in I} \lambda_{t} a_{t} = 0\) , with \(\lambda_{\kappa} = -\sum_{t \neq \kappa} \lambda_{t}\) , in other words \(\sum_{t \in I} \lambda_{t} = 0\) .

PROPOSITION 5. For a non-empty family \((a_{t})_{t\in I}\) of points of an affine space \(E\) to be affinely free, it is necessary and sufficient that, for every index \(\kappa\in I\) , \(a_{\kappa}\) do not belong to the linear variety generated by the \(a_{t}\) of index \(\neq\kappa\) .

The proposition is obvious if I has only a single element. Otherwise, taking as origin in E one of the \(a_{t}\) of index \(\neq\kappa\) , the proposition follows from §7, no.1, Remark.

